\chapter{Conclusion}

This work establishes a complete path from bitmap-based BSI algebra to GPU-based large language model inference. The earlier CPU stage showed that bit-sliced vectors could represent numeric data compactly, support a richer algebra than simple storage, and preserve control over precision and sparsity. It also showed that CPU execution was not sufficient to make those benefits practical for large neural-network workloads. That limitation motivated the transition to CUDA.

The GPU stage preserved the BSI semantics and scaled them into a transformer inference pipeline. Static weights were converted into BSI key structures, dynamic activations were converted into batched query slices, and the dominant matrix products in OPT models were executed through same-bit CUDA dot kernels. The resulting system provides substantial compression, a functioning end-to-end inference path, and a stable operating point across \texttt{opt-125m}, \texttt{opt-1.3b}, and \texttt{opt-6.7b}.

The results make the main thesis argument defensible. BSI can be used as a native compressed representation for LLM inference on GPUs, but only when representation and kernel execution are designed together. Compression alone is not enough. Lower bit count alone is not enough. The data layout, the runtime build path, and the structure of the dot kernel determine whether the representation becomes practically useful.

At the same time, the study is clear about what remains unresolved. The stable same-bit dot path still trails dense FP16 Torch execution by a large margin on larger models. Full-scope replacement also becomes increasingly sensitive as model size grows, with feed-forward layers emerging as a major quality bottleneck. These limits do not weaken the thesis. They define the current frontier precisely.

The thesis therefore closes with a clear systems conclusion: the feasibility of BSI-native transformer inference has been established, the primary tradeoffs are now measured, and the next stage of research is sharply defined. The remaining challenge is to redesign the arithmetic and layout of the BSI dot path so that the representation keeps its compression advantage while scaling much more efficiently on GPU hardware.
